\begin{lemma}[Scheide Proposition 4.2, one-tail-pair case]
Let $(T,e_0,\varphi)$ be a balanced normal globally maximum $k$-triple in
a finite loopless multigraph $G$, and suppose every vertex of $G$ is
incident with an edge.  If $h(G)+2=t(G)$, then $G$ is elementary.
\end{lemma}
\begin{proof}
Put $p=h(G)$.  The balanced condition
\noderef{BalancedGlobalMaximumTriple} identifies $p$ with the chosen
normal height.  The global equality in \noderef{GlobalMaximumTriple},
vertex injectivity in \noderef{TreeSequence}, and
\noderef{TreeSequenceVertexSet} give
\[
 \operatorname{length}(T)+1=|V(T)|=t(G)=p+2.
\]
Thus $T$ consists of the $p$ trunk vertices
$y_0,\ldots,y_{p-1}$ followed by precisely two vertices $u_1,u_2$, with
corresponding tail edges $f_1,f_2$.

The first tail edge has index $p-1\geq1$.  The Tashkinov condition contained
in \noderef{GlobalMaximumTriple} and defined in
\noderef{TashkinovTree} therefore supplies a color $\gamma$ on $f_1$ which
is missing at an earlier vertex.  Every earlier vertex lies in the trunk.
By \noderef{ScheideGraphHeightOdd}, $p$ is odd.  Taking
$j=(p+1)/2$ in the paired-tail clause of
\noderef{BalancedGlobalMaximumTriple} is legitimate because
$p<2j=p+1<p+2$; that clause gives
$\varphi(f_2)=\varphi(f_1)=\gamma$.  Hence $\gamma$ is missing on the
trunk and both tail edges have color $\gamma$.

We verify carefully that $\gamma$ is not either trunk color.  Global
maximality makes $T$ maximal for its fixed pair: a fixed-pair Tashkinov
tree strictly containing $V(T)$ would have order larger than $t(G)$, by
\noderef{TashkinovOrder} and \noderef{GlobalMaximumTriple}.  Apply
\noderef{ScheideMaximalKempePath} to this fixed-pair maximum, to the
Tashkinov tree $T$, to its first two vertices, and to the colors
$\alpha,\beta$ recorded in \noderef{NormalGlobalMaximumTriple}.  It says
that the full $(\alpha,\beta)$ component of those vertices is a path.
The normal closing path is another simple path in that same component with
the same ends.  A path has a unique simple path between its ends, so the
normal closing path is the full component.  Its recorded vertex support is
exactly $\{y_0,\ldots,y_{p-1}\}$.

The first tail edge attaches $u_1$ to an earlier, hence trunk, vertex by
\noderef{TreeSequence}.  If $\gamma$ were $\alpha$ or $\beta$, that edge
and $u_1$ would belong to the full $(\alpha,\beta)$ component, contrary to
the preceding support equality and vertex injectivity.  Thus
$\gamma\notin\{\alpha,\beta\}$.  The second tail edge initially attaches
$u_2$ either to a trunk vertex or to $u_1$.  The latter is impossible:
$f_1\ne f_2$ by edge injectivity, and two distinct $\gamma$-colored edges
incident with $u_1$ contradict the properness in
\noderef{GlobalMaximumTriple}.  Consequently both tail edges attach their
new vertices directly to the trunk.

Every non-first trunk edge is an edge-node of the normal
$(\alpha,\beta)$ closing path and is therefore colored $\alpha$ or
$\beta$.  Since $\gamma$ differs from both, all earlier non-first trunk
edges avoid $\gamma$.  We have now established the ordered block data,
equal block color, attachment-to-base condition, missing-on-base condition,
and earlier-edge avoidance required by
\noderef{ScheideStructuralConfiguration} with $s=1$, base equal to the
trunk, and the singleton color list $\Gamma=\{\gamma\}$.  The singleton
list is injective automatically.

It remains to verify the exit-covering field.  By
\noderef{ScheideGlobalMaximumOrder}, $V(T)$ is elementary.  Therefore each
of $\alpha,\beta,\gamma$ is missing at a unique tree vertex; denote a
vertex missing $\gamma$ by $y_r$.  The used tree colors are contained in
$\{\alpha,\beta,\gamma\}$, as just proved.  Every vertex other than $y_r$
misses a free color.  Indeed, a vertex other than the two ends of $e_0$
misses at least one color by
\noderef{ScheideGlobalMaximumMissingCounts}; unless it is $y_r$, none of
the three used colors can be missing there, because their missing
occurrences are already prescribed.  At either end of $e_0$, the same
counting lemma supplies at least two missing colors.  If that end is not
$y_r$, at most its own one of $\alpha,\beta$ can be among those colors, so
another missing color lies outside the three used colors.  In every case
the color obtained is missing on $V(T)$ and absent from
\noderef{TreeUsedColors}, hence belongs to
\noderef{TreeFreeColors}.

Now let $v\in\noderef{ScheideExitVertices}(T,\varphi)$.  By the definition
\noderef{ScheideExitVertex}, $v\in V(T)$.  The theorem
\noderef{ScheideExitNoFree} says that $v$ misses no free color.  The
preceding paragraph forces $v=y_r$, and therefore $\gamma$ is missing at
$v$.  This proves that every exit is covered by the sole color in
$\Gamma$.

Finally, the stated hypothesis supplies the no-isolated-vertex field, and
the original balanced normal triple supplies the globally maximum field.
Thus all fields, including the numerical length identity, of
\noderef{ScheideStructuralConfiguration} have been proved for $s=1$.
Applying \noderef{ScheideStructuralLemma} to this exact configuration
gives that $G$ is elementary.
\end{proof}
